\lysh{36885}{2}

\begin{alist}
\item Η συνάρτηση $f$ ορίζεται για τις τιμές του $x$ για τις οποίες ισχύει: $x - 2 \ne 0 \Leftrightarrow x \ne 2$.
Άρα το πεδίο ορισμού της $f$ είναι $\mathbb{A} = \mathbb{R} - \{2\}$.

\item
\begin{rlist}
\item Έχουμε ισοδύναμα:
\begin{align*}
f(x) = 0 \ &\Leftrightarrow 
\dfrac{x^2 - 1}{x - 2} = 0 \ \Leftrightarrow \\
x^2 - 1 = 0 \ &\Leftrightarrow 
x^2 = 1 \ \Leftrightarrow
\end{align*}
$x = -1 \text{ ή } x = 1$, οι οποίες είναι δεκτές.
Άρα για $x = -1$ και $x = 1, f(x) = 0$.

\item Έχουμε:
\begin{gather*}
f(0) = \dfrac{0^2 - 1}{0 - 2} = \dfrac{-1}{-2} = \dfrac{1}{2}\text{ και} \\
f(3) = \dfrac{3^2 - 1}{3 - 2} = \dfrac{8}{1} = 8.
\end{gather*}
\end{rlist}

\item Από το $\beta$) $\text{i.}$ ερώτημα έχουμε $f(-1) = 0$ και $f(1) = 0$. Από το $\beta$) $\text{ii.}$ ερώτημα έχουμε $f(0) = \dfrac{1}{2}$.
Άρα η γραφική παράσταση της $f$ τέμνει τον $x^\prime x$ άξονα στα σημεία $(-1, 0)$ και $(1, 0)$ και τον $y^\prime y$ άξονα στο σημείο $\left(0, \dfrac{1}{2}\right)$.
\end{alist}

